% Compiling Bar-Oriented Trading DSLs to Numba
% PYNE arXiv paper series — paper02-compiler
% arXiv subject: cs.PL (Programming Languages)
\documentclass[11pt,a4paper]{article}

\usepackage[margin=1in]{geometry}
\usepackage[T1]{fontenc}
\usepackage[utf8]{inputenc}

% Shared macros for PYNE arXiv paper series
% Include with: % Shared macros for PYNE arXiv paper series
% Include with: % Shared macros for PYNE arXiv paper series
% Include with: \input{../common/macros.tex}

\usepackage{amsmath,amssymb,amsfonts}
\usepackage{amsthm}
\usepackage{booktabs}
\usepackage{graphicx}
\usepackage{xcolor}
\usepackage{hyperref}
\usepackage{url}
\usepackage{algorithm}
\usepackage{algpseudocode}
\usepackage{listings}
\usepackage{tikz}
\usepackage{pgfplots}
\pgfplotsset{compat=1.17}
\usetikzlibrary{arrows.meta,positioning,fit,calc,shapes.geometric,shapes.multipart,decorations.pathreplacing,backgrounds,chains,matrix,patterns}
\usepackage{subcaption}
\usepackage{multirow}
\usepackage{tabularx}
\usepackage{enumitem}
\usepackage{microtype}
\usepackage{balance}
\usepackage{csquotes}
\usepackage{natbib}

% Colors
\definecolor{pyneorange}{HTML}{F97316}
\definecolor{pyneblue}{HTML}{2563EB}
\definecolor{pynegray}{HTML}{6B7280}
\definecolor{pynegreen}{HTML}{16A34A}
\definecolor{codebg}{HTML}{F8FAFC}
\definecolor{codekw}{HTML}{7C3AED}
\definecolor{codestr}{HTML}{B45309}
\definecolor{codecomment}{HTML}{64748B}

\hypersetup{
  colorlinks=true,
  linkcolor=pyneblue,
  citecolor=pynegreen,
  urlcolor=pyneorange,
  pdftitle={},
  pdfauthor={HOOX / PYNE Project}
}

% Code listings
\lstdefinestyle{pine}{
  basicstyle=\ttfamily\footnotesize,
  backgroundcolor=\color{codebg},
  keywordstyle=\color{codekw}\bfseries,
  stringstyle=\color{codestr},
  commentstyle=\color{codecomment}\itshape,
  breaklines=true,
  frame=single,
  rulecolor=\color{black!20},
  numbers=left,
  numberstyle=\tiny\color{pynegray},
  numbersep=6pt,
  xleftmargin=1.5em,
  framexleftmargin=1.2em,
  showstringspaces=false,
  tabsize=2,
  morekeywords={indicator,strategy,library,plot,plotshape,fill,hline,var,varip,if,else,for,while,switch,true,false,na,and,or,not,import,export,type,method,enum},
  morecomment=[l]{//},
  morestring=[b]",
}
\lstdefinestyle{python}{
  language=Python,
  basicstyle=\ttfamily\footnotesize,
  backgroundcolor=\color{codebg},
  keywordstyle=\color{codekw}\bfseries,
  stringstyle=\color{codestr},
  commentstyle=\color{codecomment}\itshape,
  breaklines=true,
  frame=single,
  rulecolor=\color{black!20},
  numbers=left,
  numberstyle=\tiny\color{pynegray},
  numbersep=6pt,
  xleftmargin=1.5em,
  framexleftmargin=1.2em,
  showstringspaces=false,
  tabsize=4,
}
\lstset{style=python}

% Theorem environments
\theoremstyle{plain}
\newtheorem{theorem}{Theorem}[section]
\newtheorem{lemma}[theorem]{Lemma}
\newtheorem{proposition}[theorem]{Proposition}
\newtheorem{corollary}[theorem]{Corollary}
\theoremstyle{definition}
\newtheorem{definition}[theorem]{Definition}
\newtheorem{example}[theorem]{Example}
\newtheorem{remark}[theorem]{Remark}

% Shortcuts
\newcommand{\pyne}{\textsc{Pyne}}
\newcommand{\pine}{Pine Script\texttrademark}
\newcommand{\tv}{TradingView\textregistered}
\newcommand{\asdl}{\textsc{Asdl}}
\newcommand{\antlr}{\textsc{Antlr}4}
\newcommand{\numba}{\textsc{Numba}}
\newcommand{\lsp}{\textsc{Lsp}}
\newcommand{\ohlcv}{\textsc{Ohlcv}}
\newcommand{\udt}{\textsc{Udt}}
\newcommand{\jit}{\textsc{Jit}}
\newcommand{\api}{\textsc{Api}}
\newcommand{\cli}{\textsc{Cli}}
% AST acronym (text mode). Keep amsmath's math asterisk \ast intact.
\DeclareRobustCommand{\Ast}{\textsc{Ast}}
% Papers in this series write \ast{} for the AST acronym in prose.
% Redefine robustly; use \mathast for a binary asterisk in formulas.
\let\mathast\ast
\DeclareRobustCommand{\ast}{\textsc{Ast}}
\newcommand{\ir}{\textsc{Ir}}
\newcommand{\ta}{\texttt{ta}}
\newcommand{\code}[1]{\texttt{#1}}
\newcommand{\file}[1]{\texttt{#1}}
\newcommand{\eg}{e.g.}
\newcommand{\ie}{i.e.}
\newcommand{\etal}{\textit{et~al.}}
\newcommand{\resp}{respectively}

% Math helpers
\newcommand{\R}{\mathbb{R}}
\newcommand{\N}{\mathbb{N}}
\newcommand{\Z}{\mathbb{Z}}
% Text-mode NA token for prose; use $\namath$ inside formulas if needed.
\DeclareRobustCommand{\na}{\textsf{na}}
\newcommand{\namath}{\mathsf{na}}
\newcommand{\series}[1]{\mathbf{#1}}
\newcommand{\baridx}{i}
\newcommand{\bars}{n}
\newcommand{\period}{p}
\newcommand{\abserr}{\varepsilon_{\mathrm{abs}}}
\newcommand{\relerr}{\varepsilon_{\mathrm{rel}}}
\newcommand{\maxerr}{\varepsilon_{\mathrm{max}}}
\newcommand{\meanerr}{\varepsilon_{\mathrm{mean}}}

% Pipeline notation — use inside math: $\pipeline$
\newcommand{\pipeline}{\mathrm{src}\!\to\!\mathrm{lex}\!\to\!\mathrm{parse}\!\to\!\mathrm{AST}\!\to\!\mathrm{eval}}

% Disclaimer footnote helper
\newcommand{\trademarknote}{%
  \pine{} and \tv{} are trademarks of TradingView, Inc.
  \pyne{} is an independent, unofficial implementation and is not affiliated with,
  authorized by, sponsored by, or endorsed by TradingView, Inc.
}

% Figure path helper (papers live one level under .papers/)
\graphicspath{{figures/}{../common/figures/}}


\usepackage{amsmath,amssymb,amsfonts}
\usepackage{amsthm}
\usepackage{booktabs}
\usepackage{graphicx}
\usepackage{xcolor}
\usepackage{hyperref}
\usepackage{url}
\usepackage{algorithm}
\usepackage{algpseudocode}
\usepackage{listings}
\usepackage{tikz}
\usepackage{pgfplots}
\pgfplotsset{compat=1.17}
\usetikzlibrary{arrows.meta,positioning,fit,calc,shapes.geometric,shapes.multipart,decorations.pathreplacing,backgrounds,chains,matrix,patterns}
\usepackage{subcaption}
\usepackage{multirow}
\usepackage{tabularx}
\usepackage{enumitem}
\usepackage{microtype}
\usepackage{balance}
\usepackage{csquotes}
\usepackage{natbib}

% Colors
\definecolor{pyneorange}{HTML}{F97316}
\definecolor{pyneblue}{HTML}{2563EB}
\definecolor{pynegray}{HTML}{6B7280}
\definecolor{pynegreen}{HTML}{16A34A}
\definecolor{codebg}{HTML}{F8FAFC}
\definecolor{codekw}{HTML}{7C3AED}
\definecolor{codestr}{HTML}{B45309}
\definecolor{codecomment}{HTML}{64748B}

\hypersetup{
  colorlinks=true,
  linkcolor=pyneblue,
  citecolor=pynegreen,
  urlcolor=pyneorange,
  pdftitle={},
  pdfauthor={HOOX / PYNE Project}
}

% Code listings
\lstdefinestyle{pine}{
  basicstyle=\ttfamily\footnotesize,
  backgroundcolor=\color{codebg},
  keywordstyle=\color{codekw}\bfseries,
  stringstyle=\color{codestr},
  commentstyle=\color{codecomment}\itshape,
  breaklines=true,
  frame=single,
  rulecolor=\color{black!20},
  numbers=left,
  numberstyle=\tiny\color{pynegray},
  numbersep=6pt,
  xleftmargin=1.5em,
  framexleftmargin=1.2em,
  showstringspaces=false,
  tabsize=2,
  morekeywords={indicator,strategy,library,plot,plotshape,fill,hline,var,varip,if,else,for,while,switch,true,false,na,and,or,not,import,export,type,method,enum},
  morecomment=[l]{//},
  morestring=[b]",
}
\lstdefinestyle{python}{
  language=Python,
  basicstyle=\ttfamily\footnotesize,
  backgroundcolor=\color{codebg},
  keywordstyle=\color{codekw}\bfseries,
  stringstyle=\color{codestr},
  commentstyle=\color{codecomment}\itshape,
  breaklines=true,
  frame=single,
  rulecolor=\color{black!20},
  numbers=left,
  numberstyle=\tiny\color{pynegray},
  numbersep=6pt,
  xleftmargin=1.5em,
  framexleftmargin=1.2em,
  showstringspaces=false,
  tabsize=4,
}
\lstset{style=python}

% Theorem environments
\theoremstyle{plain}
\newtheorem{theorem}{Theorem}[section]
\newtheorem{lemma}[theorem]{Lemma}
\newtheorem{proposition}[theorem]{Proposition}
\newtheorem{corollary}[theorem]{Corollary}
\theoremstyle{definition}
\newtheorem{definition}[theorem]{Definition}
\newtheorem{example}[theorem]{Example}
\newtheorem{remark}[theorem]{Remark}

% Shortcuts
\newcommand{\pyne}{\textsc{Pyne}}
\newcommand{\pine}{Pine Script\texttrademark}
\newcommand{\tv}{TradingView\textregistered}
\newcommand{\asdl}{\textsc{Asdl}}
\newcommand{\antlr}{\textsc{Antlr}4}
\newcommand{\numba}{\textsc{Numba}}
\newcommand{\lsp}{\textsc{Lsp}}
\newcommand{\ohlcv}{\textsc{Ohlcv}}
\newcommand{\udt}{\textsc{Udt}}
\newcommand{\jit}{\textsc{Jit}}
\newcommand{\api}{\textsc{Api}}
\newcommand{\cli}{\textsc{Cli}}
% AST acronym (text mode). Keep amsmath's math asterisk \ast intact.
\DeclareRobustCommand{\Ast}{\textsc{Ast}}
% Papers in this series write \ast{} for the AST acronym in prose.
% Redefine robustly; use \mathast for a binary asterisk in formulas.
\let\mathast\ast
\DeclareRobustCommand{\ast}{\textsc{Ast}}
\newcommand{\ir}{\textsc{Ir}}
\newcommand{\ta}{\texttt{ta}}
\newcommand{\code}[1]{\texttt{#1}}
\newcommand{\file}[1]{\texttt{#1}}
\newcommand{\eg}{e.g.}
\newcommand{\ie}{i.e.}
\newcommand{\etal}{\textit{et~al.}}
\newcommand{\resp}{respectively}

% Math helpers
\newcommand{\R}{\mathbb{R}}
\newcommand{\N}{\mathbb{N}}
\newcommand{\Z}{\mathbb{Z}}
% Text-mode NA token for prose; use $\namath$ inside formulas if needed.
\DeclareRobustCommand{\na}{\textsf{na}}
\newcommand{\namath}{\mathsf{na}}
\newcommand{\series}[1]{\mathbf{#1}}
\newcommand{\baridx}{i}
\newcommand{\bars}{n}
\newcommand{\period}{p}
\newcommand{\abserr}{\varepsilon_{\mathrm{abs}}}
\newcommand{\relerr}{\varepsilon_{\mathrm{rel}}}
\newcommand{\maxerr}{\varepsilon_{\mathrm{max}}}
\newcommand{\meanerr}{\varepsilon_{\mathrm{mean}}}

% Pipeline notation — use inside math: $\pipeline$
\newcommand{\pipeline}{\mathrm{src}\!\to\!\mathrm{lex}\!\to\!\mathrm{parse}\!\to\!\mathrm{AST}\!\to\!\mathrm{eval}}

% Disclaimer footnote helper
\newcommand{\trademarknote}{%
  \pine{} and \tv{} are trademarks of TradingView, Inc.
  \pyne{} is an independent, unofficial implementation and is not affiliated with,
  authorized by, sponsored by, or endorsed by TradingView, Inc.
}

% Figure path helper (papers live one level under .papers/)
\graphicspath{{figures/}{../common/figures/}}


\usepackage{amsmath,amssymb,amsfonts}
\usepackage{amsthm}
\usepackage{booktabs}
\usepackage{graphicx}
\usepackage{xcolor}
\usepackage{hyperref}
\usepackage{url}
\usepackage{algorithm}
\usepackage{algpseudocode}
\usepackage{listings}
\usepackage{tikz}
\usepackage{pgfplots}
\pgfplotsset{compat=1.17}
\usetikzlibrary{arrows.meta,positioning,fit,calc,shapes.geometric,shapes.multipart,decorations.pathreplacing,backgrounds,chains,matrix,patterns}
\usepackage{subcaption}
\usepackage{multirow}
\usepackage{tabularx}
\usepackage{enumitem}
\usepackage{microtype}
\usepackage{balance}
\usepackage{csquotes}
\usepackage{natbib}

% Colors
\definecolor{pyneorange}{HTML}{F97316}
\definecolor{pyneblue}{HTML}{2563EB}
\definecolor{pynegray}{HTML}{6B7280}
\definecolor{pynegreen}{HTML}{16A34A}
\definecolor{codebg}{HTML}{F8FAFC}
\definecolor{codekw}{HTML}{7C3AED}
\definecolor{codestr}{HTML}{B45309}
\definecolor{codecomment}{HTML}{64748B}

\hypersetup{
  colorlinks=true,
  linkcolor=pyneblue,
  citecolor=pynegreen,
  urlcolor=pyneorange,
  pdftitle={},
  pdfauthor={HOOX / PYNE Project}
}

% Code listings
\lstdefinestyle{pine}{
  basicstyle=\ttfamily\footnotesize,
  backgroundcolor=\color{codebg},
  keywordstyle=\color{codekw}\bfseries,
  stringstyle=\color{codestr},
  commentstyle=\color{codecomment}\itshape,
  breaklines=true,
  frame=single,
  rulecolor=\color{black!20},
  numbers=left,
  numberstyle=\tiny\color{pynegray},
  numbersep=6pt,
  xleftmargin=1.5em,
  framexleftmargin=1.2em,
  showstringspaces=false,
  tabsize=2,
  morekeywords={indicator,strategy,library,plot,plotshape,fill,hline,var,varip,if,else,for,while,switch,true,false,na,and,or,not,import,export,type,method,enum},
  morecomment=[l]{//},
  morestring=[b]",
}
\lstdefinestyle{python}{
  language=Python,
  basicstyle=\ttfamily\footnotesize,
  backgroundcolor=\color{codebg},
  keywordstyle=\color{codekw}\bfseries,
  stringstyle=\color{codestr},
  commentstyle=\color{codecomment}\itshape,
  breaklines=true,
  frame=single,
  rulecolor=\color{black!20},
  numbers=left,
  numberstyle=\tiny\color{pynegray},
  numbersep=6pt,
  xleftmargin=1.5em,
  framexleftmargin=1.2em,
  showstringspaces=false,
  tabsize=4,
}
\lstset{style=python}

% Theorem environments
\theoremstyle{plain}
\newtheorem{theorem}{Theorem}[section]
\newtheorem{lemma}[theorem]{Lemma}
\newtheorem{proposition}[theorem]{Proposition}
\newtheorem{corollary}[theorem]{Corollary}
\theoremstyle{definition}
\newtheorem{definition}[theorem]{Definition}
\newtheorem{example}[theorem]{Example}
\newtheorem{remark}[theorem]{Remark}

% Shortcuts
\newcommand{\pyne}{\textsc{Pyne}}
\newcommand{\pine}{Pine Script\texttrademark}
\newcommand{\tv}{TradingView\textregistered}
\newcommand{\asdl}{\textsc{Asdl}}
\newcommand{\antlr}{\textsc{Antlr}4}
\newcommand{\numba}{\textsc{Numba}}
\newcommand{\lsp}{\textsc{Lsp}}
\newcommand{\ohlcv}{\textsc{Ohlcv}}
\newcommand{\udt}{\textsc{Udt}}
\newcommand{\jit}{\textsc{Jit}}
\newcommand{\api}{\textsc{Api}}
\newcommand{\cli}{\textsc{Cli}}
% AST acronym (text mode). Keep amsmath's math asterisk \ast intact.
\DeclareRobustCommand{\Ast}{\textsc{Ast}}
% Papers in this series write \ast{} for the AST acronym in prose.
% Redefine robustly; use \mathast for a binary asterisk in formulas.
\let\mathast\ast
\DeclareRobustCommand{\ast}{\textsc{Ast}}
\newcommand{\ir}{\textsc{Ir}}
\newcommand{\ta}{\texttt{ta}}
\newcommand{\code}[1]{\texttt{#1}}
\newcommand{\file}[1]{\texttt{#1}}
\newcommand{\eg}{e.g.}
\newcommand{\ie}{i.e.}
\newcommand{\etal}{\textit{et~al.}}
\newcommand{\resp}{respectively}

% Math helpers
\newcommand{\R}{\mathbb{R}}
\newcommand{\N}{\mathbb{N}}
\newcommand{\Z}{\mathbb{Z}}
% Text-mode NA token for prose; use $\namath$ inside formulas if needed.
\DeclareRobustCommand{\na}{\textsf{na}}
\newcommand{\namath}{\mathsf{na}}
\newcommand{\series}[1]{\mathbf{#1}}
\newcommand{\baridx}{i}
\newcommand{\bars}{n}
\newcommand{\period}{p}
\newcommand{\abserr}{\varepsilon_{\mathrm{abs}}}
\newcommand{\relerr}{\varepsilon_{\mathrm{rel}}}
\newcommand{\maxerr}{\varepsilon_{\mathrm{max}}}
\newcommand{\meanerr}{\varepsilon_{\mathrm{mean}}}

% Pipeline notation — use inside math: $\pipeline$
\newcommand{\pipeline}{\mathrm{src}\!\to\!\mathrm{lex}\!\to\!\mathrm{parse}\!\to\!\mathrm{AST}\!\to\!\mathrm{eval}}

% Disclaimer footnote helper
\newcommand{\trademarknote}{%
  \pine{} and \tv{} are trademarks of TradingView, Inc.
  \pyne{} is an independent, unofficial implementation and is not affiliated with,
  authorized by, sponsored by, or endorsed by TradingView, Inc.
}

% Figure path helper (papers live one level under .papers/)
\graphicspath{{figures/}{../common/figures/}}

% Paper-local macros (shared set uses math-only \na)
\newcommand{\llvm}{\textsc{Llvm}}
\renewcommand{\na}{\textsf{na}}
\providecommand{\llbracket}{[\![}
\providecommand{\rrbracket}{]\!]}

\title{%
  Compiling Bar-Oriented Trading DSLs to \numba{}:\\
  Source-to-Source Lowering, Object-Mode Fallback,\\
  and Warm \ir{} Caches%
}

\author{%
  jango\_blockchained \and \pyne{} Contributors\\
  \small HOOX Research / Independent Open Source Research\\
  \small \texttt{contact@hoox.sh}%
}

\date{}

\begin{document}
\maketitle

\begin{abstract}
Bar-oriented trading domain-specific languages (\textsc{Dsl}s) evaluate
indicators and strategies by traversing a time-ordered sequence of
candlesticks. Pure interpretive execution of such programs incurs per-node
dispatch cost on every bar, which dominates wall-clock time for multi-thousand
bar histories and multi-indicator scripts. We present the \pyne{} \numba{}
compiler path: a source-to-source lowering from a Pine-compatible \ast{} into
explicit-index Python that operates on contiguous NumPy arrays and, when the
script is purely numeric, is \jit{}-compiled with \numba{}'s nopython mode to
\llvm{} machine code.

The architecture comprises four layers: (i)~a \code{CompilerVisitor} that
transpiles the \ast{} into a bar loop over \code{\_\_bar\_idx};
(ii)~a library of \code{@njit} technical-analysis kernels with incremental
$O(1)$ or amortized state updates; (iii)~flat array storage replacing
interpreter deques; and (iv)~an engine fa\c{c}ade integrating transpile,
compile, run, in-process/disk \ir{} caches, and host auto-mode with interpret
fallback. Scripts that use user-defined types, maps, or drawings are compiled
into an object-mode pure-Python bar loop that preserves those constructs while
still avoiding full \ast{} interpretation.

Internal measurements (2026) show cold-to-warm compile transitions of
$\sim$\,0.01\,ms on cache hits, warm execute medians of roughly
0.04--0.21\,ms for representative indicators at 5{,}000 bars, and
$\sim$\,49--110$\times$ end-to-end speedups versus list-based interpretation
for multi-indicator workloads after incremental-kernel fixes. Dual-host
nan-aware parity harnesses keep kernel error within floating-point noise
($\maxerr \sim 10^{-11}$). We do not claim TradingView certification or
invention of foreign multi-asset series; foreign \code{request.security} is
intentionally lowered to \na{} under a documented policy.
\end{abstract}

\paragraph{Keywords.}
source-to-source compilation, \numba{}, \jit{}, domain-specific languages,
financial indicators, dual-mode compilation, warm \ir{} caches, technical
analysis.

\paragraph{Disclaimer.}
\trademarknote{}

%==============================================================================
\section{Introduction}
\label{sec:intro}
%==============================================================================

Technical analysis systems process market data as ordered sequences of bars
(candlesticks) carrying open, high, low, close, volume, and time
(\ohlcv{}+\code{time})~\citep{murphy1999technical,tsay2010analysis}. Domain
languages in this niche---of which \pine{} is a widely used commercial
instance~\citep{tradingview-pine}---expose a declarative programming model
in which series expressions, history offsets, and sliding-window indicators
are evaluated bar-by-bar. An \ast{}-walking interpreter is a natural first
implementation: it reuses a single semantic definition for interactive
tooling, strategy simulation, and host \api{}s. Unfortunately, visitor
dispatch, dynamic attribute resolution, and list- or deque-backed series
history make the constant factors of interpretation large relative to the
arithmetic of the indicators themselves.

\pyne{} is an independent open toolchain for Pine-compatible
scripts~\citep{pynescript2026,hoox2026}. Historically it executed scripts via
a pure-Python evaluator over an \asdl{}-\ast{} produced by an \antlr{}
parser~\citep{parr2013antlr,wang1997asdl}. To approach realtime-stream and
large-backtest throughput without abandoning the Python packaging surface, the
project introduces a \emph{source-to-source} compilation path that lowers the
same \ast{} into explicit-index NumPy code and optionally \jit{}-compiles the
result with \numba{}~\citep{lam2015numba} on top of \llvm{}~\citep{lattner2004llvm}.

\paragraph{Problem.}
Let $\bars$ be the number of bars and $C$ the cost of one interpretive
expression evaluation. A script with $E$ expressions per bar has rough cost
$\Theta(\bars\cdot E\cdot C_{\mathrm{interp}})$. When $C_{\mathrm{interp}}$ is
dominated by Python bytecode and visitor overhead rather than arithmetic,
even $O(\bars)$ algorithmic complexity appears slow in wall time. Naive
window recomputation of moving averages further inflates work to
$\Theta(\bars\cdot\period)$ or worse.

\paragraph{Approach.}
We lower each script to a single closed bar loop over contiguous
\code{float64} arrays~\citep{harris2020numpy}, call precompiled \numba{}
kernels for \ta{} primitives, and retain inspectable generated Python as an
intermediate representation (\ir{}). When the \ast{} contains constructs that
cannot be represented in nopython mode (user-defined types, maps, drawings),
the same visitor emits an object-mode bar loop---still compiled away from the
\ast{} walker, but executed by CPython. An auto host mode tries compilation
first and falls back to interpretation on hard failure or known gaps, so
correctness remains primary.

\paragraph{Contributions.}
\begin{enumerate}[leftmargin=*]
  \item A four-layer compilation architecture for bar-oriented trading
        \textsc{Dsl}s: visitor lowering, numeric kernel library, flat array
        storage, and engine/runtime fa\c{c}ade (Section~\ref{sec:arch}).
  \item Detailed \ast{} lowering rules for series assignment, \code{var}/
        \code{varip}, history access, control flow, and plot allocation
        (Section~\ref{sec:lowering}).
  \item An incremental numeric kernel library with formal recurrences for
        \textsc{Sma}/\textsc{Ema}/\textsc{Rsi}/Bollinger and complexity
        analysis $O(\bars\cdot\period)$ vs.\ $O(\bars)$
        (Section~\ref{sec:kernels}).
  \item Dual-mode (numeric / object) compilation with an explicit decision
        procedure (Section~\ref{sec:object}).
  \item Multi-tier warm \ir{} caches, process prewarm, and corrupt Numba
        cache recovery (Section~\ref{sec:cache}).
  \item Host integration with \code{auto}/\code{compile}/\code{interpret}
        modes and dual-host interpret$\leftrightarrow$compile parity harnesses
        (Sections~\ref{sec:host}--\ref{sec:parity}).
  \item Evaluation of latency, throughput, and numerical residual error on
        internal 2026 benchmarks (Section~\ref{sec:eval}).
\end{enumerate}

The remainder of the paper is organized as follows.
Section~\ref{sec:bg} reviews the bar model, \numba{} modes, and \llvm{} \jit{}.
Sections~\ref{sec:arch}--\ref{sec:host} develop the design.
Section~\ref{sec:eval} reports measurements.
Sections~\ref{sec:related}--\ref{sec:concl} discuss related work, limitations,
and conclusions.

%==============================================================================
\section{Background}
\label{sec:bg}
%==============================================================================

\subsection{Bar-oriented evaluation model}
\label{sec:bg-bars}

A chart of length $\bars$ is a tuple of aligned series
\begin{equation}
  \series{O},\series{H},\series{L},\series{C},\series{V},\series{T}
  \in \R^{\bars},
\end{equation}
where $\series{T}$ stores bar-open timestamps as Unix milliseconds.
Script semantics are defined relative to the current bar index
$\baridx \in \{0,\ldots,\bars-1\}$. History access $\series{x}[k]$ denotes the
value of series $\series{x}$ at bar $\baridx-k$ (with out-of-range or missing
data mapped to the language \na{} sentinel, represented as IEEE-754
NaN~\citep{ieee754} in the numeric path).

Many indicators are sliding-window or recursive filters. A simple moving
average of period $\period$ at bar $\baridx$ is
\begin{equation}
  \label{eq:sma}
  \mathrm{SMA}_{\period}(\series{x})_{\baridx}
  =
  \begin{cases}
    \na & \text{if }\baridx < \period-1,\\[2pt]
    \dfrac{1}{\period}\displaystyle\sum_{k=0}^{\period-1}
      \series{x}_{\baridx-k}
      & \text{otherwise,}
  \end{cases}
\end{equation}
provided all window samples are finite; otherwise the result is \na{}.

Pine also provides persistent variables (\code{var}, \code{varip}) whose
initializers run once and whose values carry forward, as well as drawing and
strategy side effects that accumulate across the bar walk. These features
interact poorly with pure array vectorization when control flow or
object identity is required; an explicit bar loop remains the faithful
compilation target~\citep{nystrom2021crafting,aho2006dragon}.

\subsection{\numba{} nopython vs.\ object mode}
\label{sec:bg-numba}

\numba{} is an \llvm{}-based \jit{} for a restricted subset of
Python~\citep{lam2015numba}. In \emph{nopython} mode (\code{@njit}), the
compiler specializes functions to machine code over concrete array and scalar
types; Python objects, dictionaries of heterogeneous values, and many standard
library calls are rejected. In \emph{object} mode, \numba{} may still
accelerate loops but retains the Python object model, with substantially
weaker guarantees.

\pyne{} does not rely on \numba{} object mode for complex scripts. Instead it
implements a first-class \emph{object-mode compilation} path: pure-Python
source with a bar loop and NumPy series storage, produced by the same
visitor when nopython-incompatible constructs appear. This design choice
keeps debugging transparent (readable generated source) and avoids silent
semantic drift when \numba{} falls back internally.

\subsection{\llvm{} \jit{} and warm start}
\label{sec:bg-llvm}

\numba{} lowers specialized Python to \llvm{} \ir{} and native
code~\citep{lattner2004llvm}. First-touch compilation of a script entry
point and of the kernel library can cost hundreds of milliseconds to seconds.
Subsequent invocations reuse in-memory dispatchers; with
\code{cache=True}, Numba may also persist \code{.nbi}/\code{.nbc} artifacts.
Product systems must therefore treat cold \jit{} as a deploy-time or
readiness cost, not as interactive latency---motivating the warm-cache design
in Section~\ref{sec:cache}.

%==============================================================================
\section{Compilation Architecture}
\label{sec:arch}
%==============================================================================

Compilation is organized as four cooperating layers (Figure~\ref{fig:arch}).

\begin{figure}[t]
\centering
\begin{tikzpicture}[
  font=\small,
  box/.style={draw, rounded corners=2pt, align=center, inner sep=6pt,
              minimum width=2.6cm, minimum height=1.1cm},
  arr/.style={-{Latex}, thick},
  layer/.style={box, fill=pyneblue!8, draw=pyneblue!60},
  host/.style={box, fill=pyneorange!10, draw=pyneorange!70},
  data/.style={box, fill=pynegreen!8, draw=pynegreen!60},
]
  \node[host] (src) {Pine source\\+\ \ohlcv{}};
  \node[layer, right=0.55cm of src] (v) {Layer 1\\[1pt]
    \textbf{CompilerVisitor}\\[-1pt]
    {\footnotesize\ast{}$\to$ Python}};
  \node[layer, right=0.55cm of v] (k) {Layer 2\\[1pt]
    \textbf{numba\_builtins}\\[-1pt]
    {\footnotesize\@njit kernels}};
  \node[data, right=0.55cm of k] (a) {Layer 3\\[1pt]
    \textbf{Flat arrays}\\[-1pt]
    {\footnotesize no deques}};
  \node[host, right=0.55cm of a] (e) {Layer 4\\[1pt]
    \textbf{Engine}\\[-1pt]
    {\footnotesize transpile / run}};

  \draw[arr] (src) -- (v);
  \draw[arr] (v) -- (k);
  \draw[arr] (k) -- (a);
  \draw[arr] (a) -- (e);

  \node[below=0.85cm of v, text width=2.8cm, align=center, font=\footnotesize]
    (gen) {generated\\bar loop\\\code{\_\_bar\_idx}};
  \node[below=0.85cm of k, text width=2.8cm, align=center, font=\footnotesize]
    (lib) {\code{numba\_sma}\\\code{numba\_rsi\_inc}\\\ldots};
  \node[below=0.85cm of a, text width=2.8cm, align=center, font=\footnotesize]
    (arr) {\code{close\_arr}\\\code{my\_sma\_arr}};
  \node[below=0.85cm of e, text width=2.8cm, align=center, font=\footnotesize]
    (api) {\code{compile\_script}\\\code{Runtime.run}};

  \draw[arr, dashed, pynegray] (v) -- (gen);
  \draw[arr, dashed, pynegray] (k) -- (lib);
  \draw[arr, dashed, pynegray] (a) -- (arr);
  \draw[arr, dashed, pynegray] (e) -- (api);

  \node[draw, dashed, rounded corners, fit=(src)(v)(k)(a)(e)(gen)(lib)(arr)(api),
        inner sep=8pt, label=above:{\textbf{Four-layer \pyne{} compile pipeline}}] {};
\end{tikzpicture}
\caption{Four-layer compilation architecture. Layer~1 emits readable Python;
Layer~2 provides numeric \ta{} kernels; Layer~3 stores series as contiguous
arrays; Layer~4 packages results for hosts and manages caches.}
\label{fig:arch}
\end{figure}

\subsection{Layer 1: \code{CompilerVisitor}}

\code{CompilerVisitor} (\file{compiler.py}) is a non-evaluating \ast{} walker
in the classic visitor style~\citep{gamma1994design,aho2006dragon}. It
transpiles statements into a Python source string that indexes \ohlcv{} and
user series by an explicit bar index \code{\_\_bar\_idx}. The visitor tracks
plot titles, allocates series storage, and sets \code{object\_mode} when
non-numeric constructs appear. The generated entry point is always
\begin{lstlisting}[style=python,numbers=none]
def execute_script_compiled(open_arr, high_arr, low_arr,
                            close_arr, vol_arr, time_arr):
    ...
\end{lstlisting}
so the host contract is fixed regardless of backend.

\subsection{Layer 2: Numba builtins library}

Pine's standard library cannot be invoked from nopython code as ordinary
Python callables. Selected \ta{} and math primitives are reimplemented in
\file{numba\_builtins.py} as \code{@numba.njit(cache=True)} functions that
accept full arrays plus the current bar index (and, for incremental forms,
a small state vector \code{st}). Expanding this library is the principal path
to broader numeric coverage without leaving the \jit{} path.

\subsection{Layer 3: Flat array storage}

The compiled engine does not use interpreter \code{PineSeries} deques.
Built-in price, volume, and time series are supplied as flat NumPy arrays.
User series variables are allocated once as
\code{np.full(n\_bars, np.nan)} (or object arrays for \udt{} series) outside
the bar loop, then written at \code{\_\_bar\_idx}. This layout improves
cache locality and matches \numba{}'s preferred array
model~\citep{harris2020numpy}.

\subsection{Layer 4: Engine fa\c{c}ade}

\file{engine.py} exposes \code{transpile}, \code{compile\_script}, and
\code{run\_script}. \code{compile\_script} executes generated source in a
private namespace and returns a \code{CompiledScript} whose \code{run} method
accepts \ohlcv{} arrays. Backend \code{Runtime.run} accepts
\code{mode $\in$ \{interpret, compile, auto\}} and reshapes results into the
host response envelope (\code{plots}, \code{series}, \code{drawings},
\code{object\_mode}, cache diagnostics).

Algorithm~\ref{alg:compile} summarizes the core compile pipeline.

\begin{algorithm}[t]
\caption{\textsc{CompileScript}(source, options)}
\label{alg:compile}
\begin{algorithmic}[1]
\Require Pine source string $S$; optional flags (force object mode, cache)
\Ensure \code{CompiledScript} with \code{run} callable
\State $k_{\mathrm{raw}} \gets \mathrm{sha256}(S)$
\If{$k_{\mathrm{raw}}$ hits in-process source \textsc{Lru}}
  \State \Return cached script \Comment{warm path $\sim$0.01\,ms}
\EndIf
\State $S' \gets \mathrm{Sanitize}(S)$; $k' \gets \mathrm{sha256}(S')$
\If{$k'$ hits source \textsc{Lru} \textbf{or} disk meta for $k_{\mathrm{raw}}$}
  \State rehydrate module; share \code{execute} if \ir{} key matches
  \State \Return \code{CompiledScript}
\EndIf
\State $\mathit{AST} \gets \mathrm{Parse}(S')$
\State $(\mathit{py}, \mathit{meta}) \gets \mathrm{CompilerVisitor.visit}(\mathit{AST})$
\State $k_{\ir{}} \gets \mathrm{sha256}(\mathit{py})$
\If{$k_{\ir{}}$ hits \ir{} \textsc{Lru}}
  \State \Return script sharing warm \code{execute}
\EndIf
\State $\mathit{ns} \gets \mathrm{ExecModule}(\mathit{py})$
  \Comment{may \jit{} under \numba{}}
\State wrap nopython warm-up; on \code{TypingError} re-emit object mode
\State store source/\ir{}/disk entries; \Return \code{CompiledScript}
\end{algorithmic}
\end{algorithm}

%==============================================================================
\section{AST Lowering}
\label{sec:lowering}
%==============================================================================

This section details the principal emission rules of \code{CompilerVisitor}.
Figure~\ref{fig:lower} shows a canonical numeric example side by side.

\begin{figure}[t]
\centering
\begin{tikzpicture}[
  font=\small,
  pane/.style={draw, rounded corners=2pt, fill=codebg, inner sep=8pt,
               text width=0.42\textwidth, align=left, font=\ttfamily\scriptsize},
]
  \node[pane] (pine) {%
\textbf{\rmfamily\scriptsize Pine source}\\[4pt]
my\_sma = ta.sma(close, 14)\\
plot(my\_sma)\\
};
  \node[pane, right=0.6cm of pine] (py) {%
\textbf{\rmfamily\scriptsize Generated Python (schematic)}\\[4pt]
@numba.njit\\
def execute\_script\_compiled(...):\\
\quad n = len(close\_arr)\\
\quad my\_sma\_arr = full(n, nan)\\
\quad plot\_0 = full(n, nan)\\
\quad for \_\_bar\_idx in range(n):\\
\quad\quad my\_sma\_arr[i] = numba\_sma(...)\\
\quad\quad plot\_0[i] = my\_sma\_arr[i]\\
\quad return (plot\_0,)\\
};
  \draw[-{Latex}, very thick, pyneblue]
    (pine.east) -- node[above, font=\scriptsize\rmfamily]{lower} (py.west);
\end{tikzpicture}
\caption{Numeric-mode lowering: series assignment and \code{plot} become
full-length array writes inside an explicit bar loop. The host packs plot
titles from visitor metadata.}
\label{fig:lower}
\end{figure}

\subsection{Series assignment and storage}

A bare assignment \code{y = e} at script scope allocates \code{y\_arr} and
emits \code{y\_arr[\_\_bar\_idx] = $\llbracket e \rrbracket$} inside the loop,
where $\llbracket\cdot\rrbracket$ denotes the recursive expression
translation. Temporary expressions without names may still allocate storage
when they feed plots or later history references.

\subsection{\code{var} and \code{varip}}

Pine \code{var} initializes once on the first evaluation and carries the value
forward; \code{varip} similarly persists across realtime bar updates in the
host model. Lowering emits a guarded initializer:
\begin{lstlisting}[style=python,numbers=none]
if __bar_idx == 0:
    v_arr[__bar_idx] = <init>
else:
    v_arr[__bar_idx] = v_arr[__bar_idx - 1]
# optional subsequent assignments overwrite the carried value
\end{lstlisting}
This matches carry semantics without interpreter bookkeeping tables.

\subsection{History access}

An expression \code{close[n]} lowers to a NaN-safe index computation:
\begin{equation}
  \llbracket \series{x}[k] \rrbracket
  =
  \begin{cases}
    \series{x}_{\baridx - k} & 0 \le \baridx-k < \bars,\\
    \na & \text{otherwise,}
  \end{cases}
\end{equation}
with integer coercion that treats non-finite offsets as out-of-range.
Combinations such as \code{high[ta.highestbars(...)]} therefore stay
in-bounds when the window helper returns a negative bars-back offset
(TradingView / interpret contract).

\subsection{Control flow}

\code{if}/\code{else}, \code{for}, \code{while}, and \code{switch} lower to
the corresponding Python constructs inside the bar loop. Branching that
assigns series must write the active branch's value at \code{\_\_bar\_idx};
inactive paths leave prior NaN or carried \code{var} state depending on Pine
rules. Loops that walk bars in user code are nested inside the outer
\code{\_\_bar\_idx} traversal---a potential quadratic pattern that hosts
may later warn about, but which remains semantically faithful.

\subsection{Plot, fill, and output allocation}

Each \code{plot(expr, title=...)} allocates a sink array and assigns
$\llbracket\mathrm{expr}\rrbracket$ each bar. Titled \code{fill(plot1, plot2,
\ldots)} exports a series key for interpret/compile plot-key parity and
downstream chart band wiring; the compile path leaves the fill series as
float NaN while retaining band color / plot references in plot metadata.
Numeric mode returns a tuple of series arrays (host packs titles); object
mode returns a dictionary that may also include \code{\_\_drawings}.

\subsection{Time and \code{request.security}}

The host always supplies \code{time\_arr} (Unix ms). When the caller omits
time, the engine synthesizes \code{bar $\times$ 60\,000}. Bare \code{time},
\code{time\_close}, and calendar parts lower against \code{time\_arr}.

For \code{request.security}, compile lowers \emph{same-symbol simple \ohlcv{}
expressions} as chart series passthrough only. Foreign tickers and complex
expressions (user-defined functions, non-trivial aggregations) emit \na{}---no
invention of chart close as dividends or fundamentals. This aligns with the
interpret foreign-\na{} policy.

\subsection{Tuple unpack and inputs}

Tuple unpack such as \code{[u,m,l] = ta.bb(...)} expands to multiple array
writes from a multi-return kernel. \code{input.*} defaults are resolved at
compile/load time into constants or host-provided overrides; non-empty dynamic
input maps may gate auto-mode toward interpret for safety
(Section~\ref{sec:host}).

%==============================================================================
\section{Numeric Kernel Library}
\label{sec:kernels}
%==============================================================================

\subsection{Design constraints}

Kernels must: (i)~be nopython-legal; (ii)~accept full arrays + bar index;
(iii)~agree with the interpret oracle within floating-point noise on shared
goldens; (iv)~prefer incremental $O(1)$ amortized updates for sequential bar
walks. Full-window recomputation kernels remain available as reference
implementations and for non-sequential access.

\subsection{SMA and sliding-window sum}

Equation~\eqref{eq:sma} costs $O(\period)$ per bar if re-summed, hence
$O(\bars\cdot\period)$ per script. The incremental form maintains a running
sum $s$:
\begin{equation}
  \label{eq:sma-inc}
  s_{\baridx}
  =
  \begin{cases}
    \sum_{k=0}^{\period-1} \series{x}_{k}
      & \baridx = \period-1,\\[4pt]
    s_{\baridx-1}
      + \series{x}_{\baridx}
      - \series{x}_{\baridx-\period}
      & \baridx \ge \period,
  \end{cases}
\end{equation}
with NaN poisoning rules matching the full kernel (any non-finite window
sample yields \na{} and may force reseed). Algorithm~\ref{alg:sma}
sketches the state-vector update used by \code{numba\_sma\_inc}.

\begin{algorithm}[t]
\caption{\textsc{IncrementalSMA}($x$, $p$, $i$, $st$)}
\label{alg:sma}
\begin{algorithmic}[1]
\Require array $x$, period $p$, bar $i$, state $st=[s,\mathit{last}]$
\Ensure $\mathrm{SMA}$ at $i$; mutates $st$
\If{$i < 0$ \textbf{or} $p \le 0$} \Return \na{}
\EndIf
\State $\mathit{last} \gets$ integer from $st[1]$, or $-1$ if unset
\If{$i < \mathit{last}$} \Comment{rewind}
  \State $st[0]\gets\na{}$; $\mathit{last}\gets -1$
\EndIf
\For{$j = \mathit{last}+1$ \textbf{to} $i$}
  \If{$j < p-1$} $s \gets \na{}$
  \ElsIf{$j = p-1$} $s \gets$ sum of $x[0..p)$ if all finite else \na{}
  \Else{} update $s$ via~\eqref{eq:sma-inc}; reseed if $s$ was \na{}
  \EndIf
\EndFor
\State $st[0]\gets s$; $st[1]\gets i$
\State \Return $s / p$ if $s$ finite else \na{}
\end{algorithmic}
\end{algorithm}

\paragraph{Complexity.}
Naive full SMA over a script: $\Theta(\bars\cdot\period)$ additions.
Incremental sequential walk: $\Theta(\bars)$ additions after $O(\period)$
seed, \ie{} $O(\bars)$ total---independent of $\period$ in the steady state.

\subsection{EMA and Wilder RMA}

The exponential moving average with smoothing $\alpha = 2/(\period+1)$ is
\begin{equation}
  \label{eq:ema}
  e_{\baridx}
  =
  \alpha\,\series{x}_{\baridx}
  + (1-\alpha)\,e_{\baridx-1},
\end{equation}
after an SMA seed over the first $\period$ finite samples (compile-family
contract). Wilder's running moving average (RMA) uses
$\alpha = 1/\period$:
\begin{equation}
  \label{eq:rma}
  r_{\baridx}
  =
  \frac{1}{\period}\series{x}_{\baridx}
  + \Bigl(1-\frac{1}{\period}\Bigr) r_{\baridx-1}.
\end{equation}
Both admit $O(1)$ per-bar incremental updates with a length-2 state vector
$[\mathit{value}, \mathit{last\_i}]$ and gap catch-up when bars are skipped.

\subsection{RSI (Wilder)}

Relative Strength Index~\citep{wilder1978new} seeds average gain/loss with
the SMA of the first $\period$ price deltas, then applies RMA:
\begin{align}
  \label{eq:rsi}
  g_{\baridx} &= \max(\Delta_{\baridx}, 0),
  &
  \ell_{\baridx} &= \max(-\Delta_{\baridx}, 0),
  \\
  \bar{g}_{\baridx}
    &= \tfrac{1}{\period} g_{\baridx}
      + \bigl(1-\tfrac{1}{\period}\bigr)\bar{g}_{\baridx-1},
  \\
  \mathrm{RSI}
    &= 100 - \frac{100}{1 + \bar{g}/\bar{\ell}},
\end{align}
with the convention $\mathrm{RSI}=100$ when $\bar{\ell}=0$ after a valid
seed. The first valid bar is $\baridx=\period$ ($\period$ deltas need
$\period+1$ prices). This matches the interpret oracle and TradingView
\code{ta.rsi}; earlier mistaken ``rolling mean of gains'' implementations
were corrected in residual work (2026-08).

\subsection{Bollinger bands}

With middle band $m=\mathrm{SMA}_{\period}(x)$ and standard deviation
$\sigma$ over the same window~\citep{bollinger2001},
\begin{equation}
  \label{eq:bb}
  \mathrm{upper} = m + k\sigma,
  \qquad
  \mathrm{lower} = m - k\sigma,
\end{equation}
typically $k=2$. Incremental forms maintain the sliding sum and sum of
squares (or an equivalent stable update~\citep{higham2002accuracy}) so that
$\sigma$ is $O(1)$ amortized per bar.

\subsection{Range statistics and rank}

\code{highest}/\code{lowest} maintain window extrema; \code{highestbars}/
\code{lowestbars} return \emph{negative} bars-back offsets to the extremum,
or $-1.0$ when the window is all-\na{}. \code{ta.rci} implements Spearman
rank correlation~\citep{spearman1904} over a fixed window, matching the
interpret rank contract.

\subsection{Kernel inventory (landed subset)}

Representative nopython kernels include:
\code{sma}, \code{ema}, \code{rma}, \code{wma}, \code{hma}, \code{rsi},
\code{atr}, \code{bb}, \code{bbw}, \code{macd}, \code{stdev}, \code{cci},
\code{stoch}, \code{tsi}, \code{cmo}, \code{wpr}, \code{median},
\code{rci}, \code{vwap}, \code{linreg}, and many \code{*\_inc} variants.
Round~7 residual work landed \code{median}, \code{wpr}, \code{cmo}, and
\code{bbw} on the nopython path.

%==============================================================================
\section{Object-Mode Compilation}
\label{sec:object}
%==============================================================================

\subsection{When object mode is selected}

Figure~\ref{fig:decision} shows the decision tree. The visitor sets
\code{object\_mode} when it observes:
\begin{itemize}[leftmargin=*]
  \item user-defined type (\udt{}) definitions or instantiations;
  \item \code{map} operations (heterogeneous dictionaries);
  \item drawing \api{}s (\code{line}, \code{label}, \code{box}, \code{table},
        \ldots);
  \item strategy side effects that require broker state beyond pure series;
  \item string/color series that escape numeric representation;
  \item library \code{import} surfaces that are not kernelized;
  \item explicit \code{force\_object\_mode} after nopython warm-up
        \code{TypingError}.
\end{itemize}

\begin{figure}[t]
\centering
\begin{tikzpicture}[
  font=\small,
  dec/.style={draw, diamond, aspect=2.2, align=center, inner sep=1pt,
              fill=pyneorange!12, draw=pyneorange!70},
  act/.style={draw, rounded corners=2pt, align=center, inner sep=5pt,
              minimum width=2.4cm, fill=pyneblue!8, draw=pyneblue!60},
  term/.style={draw, rounded corners=2pt, align=center, inner sep=5pt,
              minimum width=2.4cm, fill=pynegreen!10, draw=pynegreen!60},
  arr/.style={-{Latex}, thick},
]
  \node[act] (ast) {\ast{} walk};
  \node[dec, below=0.55cm of ast] (u) {UDT / map /\\draw / strategy?};
  \node[term, below left=0.9cm and 0.1cm of u] (obj) {\textbf{Object mode}\\
    pure-Python loop\\dicts + drawings};
  \node[dec, below right=0.9cm and 0.1cm of u] (n) {Numba\\present?};
  \node[term, below left=0.85cm and -0.3cm of n] (num) {\textbf{Numeric}\\
    \code{@njit} loop};
  \node[act, below right=0.85cm and -0.3cm of n] (fb) {Object emit\\(no Numba)};
  \node[dec, below=1.5cm of num] (ty) {Warm-up\\TypingError?};
  \node[act, below left=0.7cm and 0.2cm of ty] (rec) {Re-emit\\force object};
  \node[term, below right=0.7cm and 0.2cm of ty] (ok) {Run numeric};

  \draw[arr] (ast) -- (u);
  \draw[arr] (u) -- node[left, font=\scriptsize]{yes} (obj);
  \draw[arr] (u) -- node[right, font=\scriptsize]{no} (n);
  \draw[arr] (n) -- node[left, font=\scriptsize]{yes} (num);
  \draw[arr] (n) -- node[right, font=\scriptsize]{no} (fb);
  \draw[arr] (num) -- (ty);
  \draw[arr] (ty) -- node[left, font=\scriptsize]{yes} (rec);
  \draw[arr] (ty) -- node[right, font=\scriptsize]{no} (ok);
  \draw[arr] (rec) |- (obj);
\end{tikzpicture}
\caption{Numeric vs.\ object mode decision tree. Structural features force
object emit; nopython warm-up failures recover by re-visiting with
\code{force\_object\_mode=True}.}
\label{fig:decision}
\end{figure}

\subsection{Representation}

In object mode:
\begin{itemize}[leftmargin=*]
  \item \udt{} instances are field dictionaries (and nested structures as
        nested dicts where supported);
  \item maps are ordinary Python \code{dict}s;
  \item drawing calls append structured events to a \code{\_\_drawings} list;
  \item plot series remain full-length NumPy arrays for host parity.
\end{itemize}
Arithmetic involving Pine \na{}/\code{None} is routed through \code{na\_num}
helpers. The response includes \code{object\_mode=True} so callers can
distinguish backends.

\subsection{Performance trade-off}

Object mode forgoes peak \jit{} throughput but still avoids per-expression
\ast{} visitation. For drawing-heavy or \udt{}-centric scripts it is the
correctness-preserving default; numeric kernels remain available for pure
series subexpressions when star-imported helpers are pure Python wrappers.

%==============================================================================
\section{Caching and Warm Start}
\label{sec:cache}
%==============================================================================

\subsection{Cache layers}

Product warm-compile (roadmap item H2) stacks several layers
(Table~\ref{tab:cache}):

\begin{table}[t]
\centering
\caption{Compile cache layers and scope.}
\label{tab:cache}
\begin{tabular}{@{}llll@{}}
\toprule
\textbf{Layer} & \textbf{Key} & \textbf{Bound} & \textbf{Cross-process} \\
\midrule
Source \textsc{Lru} & sha256(raw/sanitized) & 128 & no \\
\ir{} \textsc{Lru} & sha256(generated Python) & 64 & no \\
Host Runtime cache & raw source sha256 & bounded & no \\
Disk module + meta & under cache dir & on by default & yes \\
Numba \code{.nbc} & function identity & file-backed & yes \\
\bottomrule
\end{tabular}
\end{table}

Environment defaults favor production: \code{PYNE\_COMPILE\_DISK\_CACHE=1},
\code{PYNE\_COMPILE\_PREWARM=1}, and a persistent directory
(\code{PYNE\_COMPILE\_CACHE\_DIR}, Docker \code{/data/compile-cache}).

\subsection{State machine}

Figure~\ref{fig:warm} depicts the warm-cache state machine from process
start through steady state.

\begin{figure}[t]
\centering
\begin{tikzpicture}[
  font=\small,
  st/.style={draw, rounded corners=3pt, align=center, inner sep=6pt,
             minimum width=2.3cm, minimum height=1.0cm},
  arr/.style={-{Latex}, thick},
]
  \node[st, fill=pynegray!15] (cold) {Cold process\\empty \textsc{Lru}};
  \node[st, fill=pyneorange!15, right=1.2cm of cold] (pre) {Prewarm\\builtins};
  \node[st, fill=pyneblue!12, right=1.2cm of pre] (disk) {Disk \ir{}\\rehydrate};
  \node[st, fill=pynegreen!12, right=1.2cm of disk] (warm) {Warm hit\\$\sim$0.01\,ms};
  \node[st, fill=pyneorange!20, below=1.1cm of disk] (jit) {Cold \jit{}\\$\le$\,2\,s SLO};
  \node[st, fill=red!10, below=1.1cm of pre] (corr) {Corrupt\\.\texttt{nb*} };

  \draw[arr] (cold) -- (pre);
  \draw[arr] (pre) -- (disk);
  \draw[arr] (disk) -- (warm);
  \draw[arr] (disk) -- node[right, font=\scriptsize]{miss} (jit);
  \draw[arr] (jit) -- node[below, font=\scriptsize]{store} (warm);
  \draw[arr] (corr) -- node[left, font=\scriptsize]{purge+retry} (pre);
  \draw[arr, dashed] (warm) -- ++(0,-0.8) -| node[below, font=\scriptsize, pos=0.25]
    {steady-state /run} (warm);
\end{tikzpicture}
\caption{Warm-cache state machine. Deploy readiness should pay cold \jit{} via
prewarm; interactive paths prefer in-process and disk hits.}
\label{fig:warm}
\end{figure}

\subsection{Prewarm hooks}

\begin{itemize}[leftmargin=*]
  \item Python: \code{prewarm\_numba\_builtins()}, \code{prewarm\_scripts([...])},
        \code{ensure\_compile\_cache\_dir()};
  \item \cli{}: \code{pynescript prewarm [PATH\ldots]};
  \item Pro \api{}: \code{POST /compile/prewarm}; \code{GET /health} exposes a
        \code{compile} section;
  \item Lazy host prewarm on first non-test \code{/run} when enabled.
\end{itemize}

\subsection{Corrupt cache recovery}

Truncated \code{.nbi}/\code{.nbc} files raise \code{EOFError} or
\code{pickle.UnpicklingError} on load. Algorithm~\ref{alg:corrupt} matches
the engine's once-retry recovery.

\begin{algorithm}[t]
\caption{\textsc{CorruptCacheRecovery}($f$, args)}
\label{alg:corrupt}
\begin{algorithmic}[1]
\Require callable $f$ that may load Numba disk caches
\Ensure result of $f$ or re-raised non-corruption error
\State \textbf{try:} \Return $f(\textit{args})$
\State \textbf{on exception} $e$:
\If{\textbf{not} \textsc{IsNumbaCacheCorruption}($e$)}
  \State \textbf{raise} $e$
\EndIf
\State $n \gets \textsc{ClearNumbaFunctionCaches}()$
  \Comment{unlink \code{.nbi}/\code{.nbc}}
\State \textbf{log warning} with $n$ purged files
\State \Return $f(\textit{args})$ \Comment{single retry}
\end{algorithmic}
\end{algorithm}

Disk \ir{} clear (\code{clear\_disk\_compile\_cache}) remains a separate
operator action from Numba function-cache purge.

\subsection{Indicative SLOs}

Targets for operations (not hard CI gates), workstation + Numba:
\begin{itemize}[leftmargin=*]
  \item interpret \code{minimal} @ 2k bars: $\le 25$\,ms median;
  \item interpret \code{ta\_combo} @ 2k: $\le 200$\,ms median;
  \item cold compile first SMA-class script: $\le 2$\,s;
  \item warm compile (source cache hit): $\le 1$\,ms (typically
        $\sim$0.01--0.05\,ms);
  \item warm run \code{ta\_combo} @ 5k: $\le 5$\,ms (often $\sim$1\,ms);
  \item cross-process disk rehydrate: $\le\sim$60--70\% of full cold.
\end{itemize}

%==============================================================================
\section{Host Integration and Auto Mode}
\label{sec:host}
%==============================================================================

\subsection{Runtime modes}

\begin{description}[leftmargin=*,style=nextline]
  \item[\code{auto}] Default on Pro \api{} \code{/run} and \code{/run/batch}.
        Prefer compile when eligible; fall back to interpret on hard failure
        or known gaps.
  \item[\code{compile}] Compile only (error if transpile/exec fails;
        nopython$\to$object recovery still occurs inside compile).
  \item[\code{interpret}] \ast{} evaluator only.
\end{description}

Safe auto gates (no silent wrong results) include: non-empty dynamic
\code{inputs} $\to$ interpret; top-level \code{import} / \code{request.*}
$\to$ interpret; any compile error $\to$ interpret with
\code{compile\_fallback\_reason}. Response fields include
\code{auto\_backend}, \code{compile\_cached}, \code{compile\_ms}, and
optional \code{nopython\_fallback\_reason}.

\begin{algorithm}[t]
\caption{\textsc{AutoMode}(source, ohlcv, inputs)}
\label{alg:auto}
\begin{algorithmic}[1]
\Require script $S$, market data $D$, optional inputs $I$
\Ensure host result dict with \code{series}/\code{plots}
\If{$I$ non-empty \textbf{or} $S$ has top-level import/request}
  \State \Return \textsc{Interpret}($S,D,I$) with
         \code{auto\_backend=interpret}
\EndIf
\If{\textbf{not} \textsc{HasNumba}() \textbf{and} numeric-only path required}
  \State attempt object compile or fall back interpret
\EndIf
\State \textbf{try:} $R \gets \textsc{CompileAndRun}(S,D)$
\State \textbf{on success:} \Return $R$ with \code{auto\_backend=compile}
\State \textbf{on exception} $e$:
\State \Return \textsc{Interpret}($S,D,I$) with
       \code{compile\_fallback\_reason}=\textsc{Msg}($e$)
\end{algorithmic}
\end{algorithm}

\subsection{Engine \api{} surface}

Public entry points:
\begin{itemize}[leftmargin=*]
  \item \code{transpile(source) $\to$ str} --- generated Python for inspection;
  \item \code{compile\_script(source) $\to$ CompiledScript};
  \item \code{run\_script(source, ohlcv) $\to$ result dict};
  \item cache helpers: \code{clear\_compile\_cache},
        \code{clear\_disk\_compile\_cache},
        \code{clear\_numba\_function\_caches}, \code{compile\_cache\_stats}.
\end{itemize}

\subsection{Result envelope}

Both modes normalize into a host-facing structure with plot series, optional
drawings, bar count, identifiers, and backend flags. Numeric mode avoids
returning Numba \code{typed.Dict} across the boundary (tuple packing + host
title map) to eliminate expensive conversion observed in early prototypes.

%==============================================================================
\section{Correctness: Dual-Host Parity}
\label{sec:parity}
%==============================================================================

Compilation is only useful if results match the interpret oracle on the
supported surface. \pyne{} maintains dual-host checks rather than claiming
third-party certification.

\subsection{Harness}

\file{scripts/compare\_interp\_compile.py} runs the same scripts under
\code{mode=interpret} and \code{mode=compile}, compares
\code{result["series"]} with nan-aware \code{allclose}, and buckets residuals
(\code{OK}, \code{fill\_background\_only}, \code{both\_error\_same},
\code{MISMATCH}, \ldots). Always-on smoke tests live in
\file{tests/test\_interp\_compile\_parity.py}; focused foreign-security cases
in \file{tests/test\_dividend\_yield\_parity.py}.

\subsection{Landed correctness notes (2026-08)}

\begin{itemize}[leftmargin=*]
  \item \code{time\_arr} always present (host real timestamps or synthetic ms).
  \item \code{request.security}: same-symbol simple \ohlcv{} passthrough;
        foreign/complex $\to$ \na{}.
  \item \textsc{Rsi} Wilder seed then RMA (Section~\ref{sec:kernels}).
  \item \code{highestbars}/\code{lowestbars}: negative bars-back; all-\na{}
        $\to -1.0$.
  \item Titled \code{fill()} series keys for plot-key parity.
  \item Numba cache recovery from corrupt \code{.nbi}/\code{.nbc}.
  \item \code{ta.rci} Spearman rank.
\end{itemize}

\subsection{Error metrics}

For aligned series $a,b$ of length $\bars$, with mask
$m_i = \neg(\mathrm{isnan}(a_i)\lor\mathrm{isnan}(b_i))$,
\begin{equation}
  \maxerr
  =
  \max_{i:m_i} \lvert a_i - b_i \rvert,
  \qquad
  \meanerr
  =
  \frac{1}{\lvert m\rvert}\sum_{i:m_i}\lvert a_i-b_i\rvert.
\end{equation}
Kernel parity studies report $\maxerr \sim 10^{-11}$ attributable to floating
associativity, not algorithmic divergence~\citep{higham2002accuracy,ieee754}.
NaN patterns must also match (both \na{} or both finite) for a series pair to
be considered \code{OK}.

%==============================================================================
\section{Evaluation}
\label{sec:eval}
%==============================================================================

\subsection{Methodology}

Unless noted, measurements use workstation-class hardware, CPython with Numba
installed, synthetic rising or random-walk \ohlcv{}, and medians over repeated
runs after explicit cache clears or warm-ups as labeled. Numbers are
\emph{internal 2026 engineering benchmarks}, not a multi-platform public
reproducibility suite; they characterize order of magnitude and design
progress across compile performance rounds.

\subsection{Compile vs.\ execute latency}

Table~\ref{tab:exec} reports warm compiled execute medians at
$\bars=5000$ for representative indicators after incremental-kernel work.

\begin{table}[t]
\centering
\caption{Warm compile execute medians at $n=5000$ bars (internal 2026).}
\label{tab:exec}
\begin{tabular}{@{}l r r@{}}
\toprule
\textbf{Script} & \textbf{Median (ms)} & \textbf{Notes} \\
\midrule
SMA & $\sim$0.04 & single window MA \\
RSI & $\sim$0.08 & Wilder incremental \\
BB  & $\sim$0.15 & bands + stdev \\
TSI & $\sim$0.11 & double-smoothed \\
MULTI (SMA+EMA+RSI) & $\sim$0.21 & multi-plot \\
\code{ta\_combo} warm compile hit & $\sim$0.01 & cache only \\
\code{ta\_combo} warm run & $\sim$1.06 & full combo @5k \\
\bottomrule
\end{tabular}
\end{table}

Early MACD/EMA paths re-accumulated from bar~0 on every index, producing
near-quadratic execute times ($\sim$33\,ms at 5k bars). After incremental
state vectors, MACD execute fell by $\sim$110$\times$ and MULTI by
$\sim$49$\times$ relative to those broken baselines on the same machine class
(compile+exec vs.\ list-based interpret for MACD/MULTI class workloads remains
in the tens-to-hundreds$\times$ range depending on interpret plot packing).

\subsection{Cold vs.\ warm compile}

Cold compile is dominated by Numba first-touch of the generated entry and
builtins (often $\sim$0.4--2\,s for SMA-class scripts depending on process
cache state). Warm source-\textsc{Lru} hits are $\sim$0.01--0.05\,ms. Disk
\ir{} rehydrate is cheaper than full cold but not free
(Section~\ref{sec:cache} SLOs).

\subsection{Throughput charts}

Figure~\ref{fig:perf} visualizes schematic latency bars for interpret vs.\
cold compile vs.\ warm compile+run on a multi-indicator workload, using the
internal measurements above as anchors.

\begin{figure}[t]
\centering
\begin{tikzpicture}
\begin{axis}[
  ybar,
  bar width=12pt,
  width=0.95\textwidth,
  height=6.2cm,
  ylabel={Latency (ms, log scale)},
  ymode=log,
  log origin=infty,
  ymin=0.005,
  ymax=5000,
  symbolic x coords={Interpret ta\_combo@2k,
                     Cold compile SMA,
                     Warm compile hit,
                     Warm run SMA@5k,
                     Warm run MULTI@5k,
                     Warm run ta\_combo@5k},
  xtick=data,
  x tick label style={font=\scriptsize, rotate=25, anchor=east},
  nodes near coords,
  nodes near coords style={font=\tiny, /pgf/number format/fixed},
  every node near coord/.append style={font=\tiny},
  point meta=explicit symbolic,
  enlarge x limits=0.12,
  legend style={at={(0.5,1.02)}, anchor=south, legend columns=3, font=\scriptsize},
  grid=major,
  major grid style={dotted, gray!40},
]
\addplot[fill=pynegray!50, draw=black!60] coordinates {
  (Interpret ta\_combo@2k, 153) [=153]
  (Cold compile SMA, 946) [=946]
  (Warm compile hit, 0.01) [=0.01]
  (Warm run SMA@5k, 0.04) [=0.04]
  (Warm run MULTI@5k, 0.21) [=0.21]
  (Warm run ta\_combo@5k, 1.06) [=1.06]
};
\end{axis}
\end{tikzpicture}
\caption{Indicative latencies (internal 2026): interpret combo residual vs.\
cold \jit{}, warm compile cache hits, and warm numeric executes. Log scale
emphasizes the multi-order-of-magnitude gap between cold \jit{} and warm hits.}
\label{fig:perf}
\end{figure}

Figure~\ref{fig:speedup} contrasts naive $O(n\cdot p)$ vs.\ incremental
$O(n)$ kernel cost as period grows (schematic, unit work normalized).

\begin{figure}[t]
\centering
\begin{tikzpicture}
\begin{axis}[
  width=0.85\textwidth,
  height=5.5cm,
  xlabel={Window period $p$},
  ylabel={Relative work units},
  xmin=0, xmax=220,
  ymin=0, ymax=250,
  grid=major,
  major grid style={dotted, gray!40},
  legend pos=north west,
  legend style={font=\small},
]
\addplot[thick, pyneorange, domain=1:200, samples=50] {x};
\addlegendentry{naive $O(n\cdot p)/n$ $\propto p$}
\addplot[thick, pyneblue, domain=1:200, samples=2] {8};
\addlegendentry{incremental $O(1)$ amortized}
\end{axis}
\end{tikzpicture}
\caption{Asymptotic per-bar work for sliding SMA-style kernels: naive
recompute grows with $p$; incremental state remains flat after seeding.}
\label{fig:speedup}
\end{figure}

\subsection{Parity residuals}

Across kernel golden suites, compiled vs.\ full-kernel or interpret oracles
typically achieve $\maxerr \le 10^{-10}$ to $10^{-11}$ for continuous
indicators. Structural residuals (hline-only keys, fill background series)
are classified separately from true value mismatches by the dual-host
harness.

\subsection{Object mode}

Object-mode UDT/plot scripts compile in tens of milliseconds cold (no Numba
entry \jit{}) and run on the order of single-digit milliseconds at 5k bars
in microbenchmarks---slower than numeric mode but faster than full \ast{}
interpretation for comparable expression volume, with broader language
coverage.

%==============================================================================
\section{Related Work}
\label{sec:related}
%==============================================================================

\paragraph{\numba{} and Python \jit{}.}
Lam \etal{} introduce \numba{} as an \llvm{}-based \jit{} for array-oriented
Python~\citep{lam2015numba}. Our work does not extend Numba itself; it is a
domain-specific front-end that emits Numba-legal programs and a parallel
object-mode path when the domain language exceeds nopython. Nuitka and
similar whole-program compilers~\citep{nuitka} pursue different packaging
goals without bar-series semantics.

\paragraph{Financial and trading \textsc{Dsl}s.}
Commercial platforms provide proprietary compilers for bar languages; public
literature more often discusses time-series libraries
(\eg{} pandas~\citep{brooks2015python}) than closed-form compilation of
Pine-like semantics. \pyne{} targets an open, inspectable intermediate Python
\ir{} rather than an opaque native bytecode blob.

\paragraph{Partial evaluation and staging.}
Source-to-source lowering with residual bar loops resembles online partial
evaluation~\citep{jones1993partial}: market data remains dynamic while script
structure is specialized into straight-line array code. We do not currently
apply binding-time analysis formally; the visitor encodes domain knowledge
(series vs.\ scalar, plot allocation) directly.

\paragraph{Dual-mode virtual machines.}
Self-optimizing and dual-mode engines (\eg{} object vs.\ optimized paths in
dynamic language VMs~\citep{chambers1991iterative}) inspire our
numeric/object split. The difference is that both modes are produced by
ahead-of-time source generation, not by runtime tracing of an interpreter
core.

\paragraph{Compiler construction.}
Classic pipelines~\citep{aho2006dragon,cooper2011engineering} separate
front-end, \ir{}, and back-end. We reuse ANTLR+\asdl{} as the front-end,
generated Python as a human-readable \ir{}, and Numba/\llvm{} or CPython as
back-ends---trading maximal optimization for packaging simplicity and
debuggability~\citep{leroy2009formal}.

%==============================================================================
\section{Limitations and Future Work}
\label{sec:limits}
%==============================================================================

\begin{itemize}[leftmargin=*]
  \item \textbf{No TradingView certification.} Numerical and structural
        parity is measured against \pyne{}'s interpret oracle and public
        documentation, not against a licensed reference binary.
  \item \textbf{Foreign multi-asset data.} Compile does not invent foreign
        series; complex \code{request.security} yields \na{}. Full multi-ticker
        hosts remain future work.
  \item \textbf{Strategy breadth.} Complex order/\code{StrategyState} paths
        remain interpret-first; basic entry/close exist in object mode.
  \item \textbf{Drawing fidelity.} Deletes and full style parity with
        interpreter registries are incomplete.
  \item \textbf{Nested UDTs/methods.} Coverage is partial; object mode is the
        escape hatch.
  \item \textbf{True AOT.} Disk + Numba \code{.nbc} is not a portable
        ahead-of-time binary; multi-worker deploys still warm per process
        unless coordinated.
  \item \textbf{mypyc / native kernels.} Deferred; residual interpret plot
        packing still dominates some combo walls on the non-compile path.
\end{itemize}

Future work prioritizes expanding the nopython kernel surface so fewer
scripts fall out of numeric mode, unifying dual-host Runtime packages, and
growing the interpret$\leftrightarrow$compile golden corpus under the parity
harness rather than ad-hoc scripts.

%==============================================================================
\section{Conclusion}
\label{sec:concl}
%==============================================================================

We presented the \pyne{} \numba{} compiler path for bar-oriented trading
\textsc{Dsl}s: a four-layer architecture that lowers a Pine-compatible \ast{}
to explicit-index Python, executes pure numeric scripts under \numba{}
nopython \jit{}, and falls back to an object-mode bar loop for \udt{}s, maps,
and drawings. Incremental $O(\bars)$ kernels, multi-tier warm \ir{} caches,
corrupt-cache recovery, and host auto-mode with interpret fallback make the
system practical for interactive \api{}s and large histories.

Internal 2026 measurements show warm compile hits near $10^{-2}$\,ms, warm
indicator executes on the order of $10^{-1}$\,ms at 5{,}000 bars, and large
speedups versus list-based interpretation once quadratic kernel bugs were
eliminated. Dual-host nan-aware parity keeps residuals at floating-point
noise for landed kernels. The design emphasizes inspectable intermediates and
correctness-first auto fallback over unverifiable peak claims.

\section*{Acknowledgments}

We thank \pyne{} and HOOX contributors for performance rounds, parity
harnesses, and production warm-compile integration. \trademarknote{}

%==============================================================================
\section*{References}
%==============================================================================

\bibliographystyle{plainnat}
\bibliography{../common/bibliography}

\end{document}
